Second, the institutional assumption about PPA1 timing is not observed at the individual level. We know the operational first-year context and the three-visit threshold used in the study, but we do not directly observe which visit, if any, a student undertook specifically to satisfy a PPA requirement. We also do not observe internal motivation.

Third, the primary cohort conditions on passing MU and progressing to Calculus with valid linked records. This defines a substantively meaningful risk set for later Calculus use but also selects on academic progression. Results therefore generalize most directly to students who follow that pathway, not to all first-year students.

Fourth, visit records measure recorded encounters rather than duration, intensity, content quality, or learning processes. Multiple same-day visits can represent legitimate continued support across adviser blocks, and the data do not support treating them as manipulation.

Fifth, the study is based on one institution and one support-centre context. The contribution therefore depends on connecting the local institutional design to broader higher-education processes rather than assuming direct generalizability. Cross-institution replication would be valuable.

Finally, the strict $N=3{,}241$ longitudinal analysis was produced in a later development stage than the currently reproducible methodology-report pipeline. The underlying analysis and retained outputs are documented, but the longitudinal logic still needs to be fully reintegrated into a single canonical executable pipeline before submission.

\draftnote{Submission should wait until the $N=3{,}241$ analysis can be regenerated from the canonical runner on the controlled source data, with a provenance manifest linking all paper numbers and the figure to that run.}

\draftnote{If feasible, add pre-MU covariates such as admission/diagnostic mathematics measures. These would not create causal identification, but they would improve the selection model and strengthen interpretation of the adjusted persistence gradient.}

\draftnote{Consider a short survey or qualitative follow-up on reasons for first and later CMAT use. Administrative sequences can establish behavioural persistence, but mechanism claims about familiarity, motivation, or perceived usefulness require direct evidence.}

\section{Conclusion}

This study shows that formal mathematics-support use can persist well beyond an incentive-linked first-year context. Among students progressing from MU to Calculus, later CMAT use increased sharply with prior use, from 17\% among non-users to 67\% among students with four or more MU visits. Students who continued beyond the operational three-visit threshold were also more likely to return later than students who stopped exactly at three. These patterns survive adjustment for prior relative performance, degree programme, and term.

The evidence does not show that PPA caused persistence or that CMAT attendance caused academic improvement. Its contribution is instead to separate three objects that are often conflated in evaluations of voluntary support: institutional conditions that encourage initial participation, student-selected use while those conditions are active, and later continuation when the same incentive context no longer applies. Treating these as distinct outcomes provides a more informative way to study academic support, engagement, and persistence in higher education.

\bibliographystyle{unsrtnat}
\bibliography{references}

\end{document}
